\section{Conclusion}
\label{sec:conclusion}

Repository-level repair evidence is produced through a sequence of measurements. A logged action field must satisfy its operational invariant; a governed action must still reach the verifier; repeated opportunities must be separated from broader task support; and a verifier pass retains an explicit semantic boundary.

The \TotalRows-execution audit shows why these distinctions matter. Twelve raw action positives fail the application invariant used by the governed-action construct. After harmonization, conditional endpoint yield is \CleanYield\% and \ExtensionYield\% in the two DeepSeek strata. At eight retained opportunities, action discovery reaches \CleanActionDiscoveryKEight\% and \ExtensionActionDiscoveryKEight\%, while pass discovery reaches \CleanPassDiscoveryKEight\% and \ExtensionPassDiscoveryKEight\%. The 28 passing rows occupy seven cells, three tasks, and two task families. Endpoint controls confirm the expected behavior on the tested reference, no-op, and mutant cases.

The claim index binds each result to construct, opportunity design, support unit, endpoint interpretation, and instrumentation provenance. This specification keeps repair statistics comparable and makes intermediate evidence useful on its own terms. A number becomes a scientific result when the claim it measures is equally explicit.
